\chapter{克莱因曲面及其基本群讲义}
\label{chap:02}

\noindent Florent Schaffhauser\footnote{感谢所有参加者的关注与富有启发的提问。特别感谢研讨会组织者
Steven Bradlow 与 Óscar García Prada，研究计划组织者 Luis Álvarez-Cónsul 与 Ignasi Mundet i
Riera，以及报告人与参加者 Olivier Guichard、Jacques Hurtubise、Andrés Jaramillo Puentes、
Melissa Liu、Gregor Masbaum、Toni Pantev 与 Richard Wentworth。}

本讲义基于 2012 年在巴塞罗那 CRM 所作的三次一小时报告，属于「曲面群表示论大师班与
研讨会」（Master Class and Workshop on Representations of Surface Groups）活动的一部分，
而该活动本身又是「模空间的几何与量子化」研究计划的一部分。这些报告的目标是介绍克莱因
曲面的一般理论，尤其是这类曲面上「基本群」的恰当概念，并始终强调它与实代数几何中
更偏代数的基本群视角之间的类比。

确实，空间 $X$ 的拓扑基本群通常用 $X$ 中的道路来定义，即给定基点 $x\in X$ 处环路的
同伦类之集。然后在关于 $X$ 的温和拓扑假设下，构造 $(X,x)$ 的万有覆盖 $(\tilde X,\tilde x)$，
并证明 $X$ 在 $x$ 处的基本群同构于该万有覆盖的自同构群，而中间覆盖对应于基本群的
子群。在 SGA-1（重印于 [6]）中，Grothendieck 把基本群的这最后一条性质——它分类连通
空间 $X$ 的覆盖——用作域上连通概形的代数基本群的定义：此时代数基本群是 $X$ 的适当
覆盖范畴的 Galois 群，并且是基域绝对 Galois 群的一个增广（augmentation）。他还证明了，
当 $X$ 是复代数簇时，其代数基本群是相应复解析空间的拓扑基本群的 profinite 完备化。
把视角反过来，可以说黎曼面的拓扑基本群是复代数曲线的代数基本群的离散类比。从这个
观点看，本讲义的目标就是解释：当复曲线实际定义在实数上时，哪个群是代数基本群的
离散类比。

从复解析的角度看，实代数曲线是一对 $(X,\sigma)$，其中 $X$ 是黎曼面，$\sigma$ 是 $X$ 的
反全纯对合。按照 SGA-1 中所述代数基本群的性质，$(X,\sigma)$ 的离散基本群应当是一个群
$\pi^R_1((X,\sigma);x)$，它落在如下短正合列中
\[
1\longrightarrow \pi_1(X;x)\longrightarrow \pi^R_1((X,\sigma);x)
\longrightarrow \Gal(\C/\R)\longrightarrow 1,
\]
其中 $\pi_1(X;x)$ 是 $X$ 的拓扑基本群。上面的正合列说明，$\pi^R_1((X,\sigma);x)$
实际上就是 $X/\sigma$ 的轨形基本群，它最初由 Thurston 用轨形道路定义（轨形的一般理论
与 Thurston 的定义可参看 [15, 17]）。在本讲义中，我们改用实曲线 $(X,\sigma)$ 的
覆盖空间来给出 $\pi^R_1((X,\sigma);x)$ 的定义。

第 2.1 节简要叙述 Real 黎曼面与克莱因曲面的一般理论，包括紧连通 Real 黎曼面的拓扑
分类。在第 2.2 节中，我们定义克莱因曲面的基本群，并说明它如何分类 $(X,\sigma)$ 的实
覆盖空间。在第三节即最后一节 2.3 中，我们研究克莱因曲面基本群的线性表示与酉表示，
并说明轨形基本群 $\pi^R_1((X,\sigma);x)$ 的酉表示簇如何映入拓扑基本群
$\pi_1(X;x)$ 的酉表示簇，且其像包含在某个对合的不动点集之中。

\section{克莱因曲面与实代数曲线}
\label{sec:02.1}

本节回顾克莱因曲面理论中的一些熟知材料。更完整的参考文献可参看综述文章
Natanzon [14] 与 Huisman [8]，以及 Alling 与 Greenleaf 的专著 [1]。

\subsection{代数曲线与二维流形}
\label{sec:02.1.1}

自 Riemann 的工作以来，一个非常熟悉的观念是：与非奇异复代数曲线相联系的是一个可定向
（事实上带定向）的实曲面，即一个二维流形。反过来，任何紧连通可定向曲面都容许一个
一维复解析流形结构（即黎曼面结构），并且相对于该结构，它可以嵌入为 $\CP^3$ 的复子
流形。由 Chow 定理，这种嵌入的像是代数的。本小节的目标是简要回顾：实代数曲线同样有
与之相联系的曲面，而这些曲面的拓扑与几何编码了曲线的一些代数性质。

我们从一个简单的例子开始。设 $P\in\R[x,y]$ 是二元实系数非常数多项式。当然 $P$ 也可
视为 $\C[x,y]$ 的元素。令
\[
X_P(\C):=\{(u,v)\in\C^2\mid P(u,v)=0\},
\qquad
X_P(\R):=\{(u,v)\in\R^2\mid P(u,v)=0\}.
\]
注意 $X_P(\R)$ 是实平面曲线，$X_P(\C)$ 是复平面曲线。复仿射平面 $\C^2$ 带有一个实结构
（即一个反全纯对合）
\[
\sigma:\ \C^2\longrightarrow\C^2,\qquad (u,v)\longmapsto(\bar u,\bar v),
\]
而实仿射平面 $\R^2\subset\C^2$ 等于 $\Fix(\sigma)$。由于 $P$ 的系数是实的，$X_P(\C)$
是 $\sigma$-不变的，并且有 $X_P(\R)=\Fix\bigl(\sigma|_{X_P(\C)}\bigr)$。换言之，方程
$P=0$ 的实解恰好就是 $\sigma$-不变的复解。注意 $X_P(\R)$ 可以是空集；而由 Hilbert
零点定理（Nullstellensatz），$X_P(\C)$ 总是非空的。要点在于：用对子 $(X_P(\C),\sigma)$
代替 $X_P(\R)$ 有一个直接的好处——人们确信自己永远不会在谈论空集。但这个故事背后的
结构还不止于此。

设偏导数 $\partial P/\partial u$ 与 $\partial P/\partial v$ 不在 $X_P(\C)$ 的点上同时
为零。那么在 $\C^2$ 的通常拓扑所诱导的拓扑下，复平面曲线 $X_P(\C)$ 具有黎曼面结构
（即一个局部以 $\C$ 的开集为模型的复解析图册；例如可参看 [4]），而 $\C^2$ 的实结构
$\sigma:(u,v)\mapsto(\bar u,\bar v)$ 诱导出 $X_P(\C)$ 的一个反全纯对合，其不动点集
正是 $X_P(\R)$。上面的例子引出如下定义。

\begin{definition}[Real 黎曼面]\label{def:2.1.1}
\emph{Real 黎曼面}是一个对子 $(\Sigma,\tau)$，其中 $\Sigma$ 是黎曼面，$\tau$ 是 $\Sigma$
的反全纯对合：
\begin{enumerate}[label=(\roman*), leftmargin=2.2em, itemsep=0.15em]
\item 对每个 $x\in\Sigma$，切映射 $T_x\tau:T_x\Sigma\to T_{\tau(x)}\Sigma$ 是
      $\C$-反线性的（特别地 $\tau\neq\Id_\Sigma$）；
\item $\tau^2=\Id_\Sigma$。
\end{enumerate}
这样的 $\tau$ 称为 $\Sigma$ 上的一个\emph{实结构}。
\end{definition}

Real 黎曼面的同态 $f:(\Sigma,\tau)\to(\Sigma',\tau')$ 是一个全纯映射
$f:\Sigma\to\Sigma'$，使得 $f\circ\tau=\tau'\circ f$。

若 $\Sigma$ 连通（相应地紧），则称 Real 黎曼面 $(\Sigma,\tau)$ 连通（相应地紧）。
Real 黎曼面是 Atiyah 意义下 Real 空间（见 [2]）的一个例子。我们之所以把 Real 的
首字母大写，正是遵循 Atiyah 的约定。下文我们也常说实代数曲线或实曲线来代替
Real 黎曼面。

\begin{annote}
「Real」大写是 Atiyah 的记号，指带反全纯对合（实结构）的复对象，与「定义在 $\R$ 上」
并不等同：例如 $\CP^1$ 上的 $\tau_0:z\mapsto-1/\bar z$ 是实结构，但没有实点。
为避免与普通形容词「实」混淆，译文保留 Real 不译。
\end{annote}

下一个例子说明：给定实曲线 $(\Sigma,\tau)$，商曲面 $\Sigma/\tau$ 具有代数意义
（我们用 $\Sigma/\tau$ 记 $\Sigma$ 对由 $\tau$ 定义的 $\Z/2\Z$-作用的商）。令
$V_\R:=\Spm\R[T]$ 为环 $\R[T]$ 的极大理想之集，并记 $V_\C:=\Spm\C[T]$。$\C[T]$
（相应地 $\R[T]$）的极大理想是由不可约多项式生成的理想，故
\[
V_\C=\{(T-z): z\in\C\},
\qquad
V_\R=\{(T-a): a\in\R\}\cup\{(T^2+bT+c): b,c\in\R,\ b^2-4c<0\}.
\]
特别地，有双射 $V_\C\simeq\C$，而 $\C$ 的实结构 $\sigma:z\mapsto\bar z$ 诱导出
$V_\C$ 的对合 $(T-z)\mapsto(T-\bar z)$，其不动点集为
$\{(T-a):a\in\R\}\simeq\R=\Fix(\sigma)\subset\C$。

若 $\operatorname{Im}z\neq 0$，则集合 $\{z,\bar z\}$ 是实多项式
$T^2-(z+\bar z)T+z\bar z$ 的根集，其判别式为
\[
b^2-4c=(z+\bar z)^2-4|z|^2=(z-\bar z)^2=-4(\operatorname{Im}z)^2<0.
\]
反之，形如 $T^2+bT+c$ 且满足 $b^2-4c<0$ 的实多项式有唯一的复根 $z$ 使
$\operatorname{Im}z>0$，另一根为 $\bar z$。换言之，有双射
\[
V_\R\simeq\bigl\{\{\Id_\C,\sigma\}\simeq\Gal(\C/\R)\ \text{在}\ V_\C\ \text{中的轨道}\bigr\}
=V_\C/\sigma .
\]
特别地，作为集合，
\[
V_\C/\sigma=\C/\sigma\simeq\{z\in\C\mid \operatorname{Im}z\ge 0\}=:\C^+
\]
是「闭」上半平面。于是，代数对象 $V_\R=\Spm\R[T]$ 自然地与一个带边曲面——即闭上半
平面 $\C/\sigma$——双射。这提示我们考虑形如 $\Sigma/\tau$ 的一般商，其中
$(\Sigma,\tau)$ 是定义 2.1.1 意义下的 Real 黎曼面。

按构造，$\Sigma/\tau$ 是带边的二维流形。它可以是不可定向的，边界可以是空的。首先
要观察的是：$\Sigma/\tau$ 上有良好的函数概念。事实上，记 $p:\Sigma\to\Sigma/\tau$ 为
典范投影，$\calO_\Sigma$ 为 $\Sigma$ 上全纯函数层，则群
$\Z/2\Z=\{\Id_\Sigma,\tau\}$ 通过
\begin{equation}\label{eq:2.1}
(\tau\cdot f)(x):=f\bigl(\tau(x)\bigr)
\end{equation}
作用在直像 $p_*\calO_\Sigma$ 上。这是良定义的，因为按定义 $p_*\calO_\Sigma$ 的局部
截面是定义在 $\Sigma$ 的某个 $\tau$-不变开子集上的全纯函数，而按上式定义的 $\tau\cdot f$
亦然。记 $(p_*\calO_\Sigma)^\tau$ 为 $p_*\calO_\Sigma$ 中 $\tau$-不变局部截面构成的
子层。于是 $(p_*\calO_\Sigma)^\tau$ 的任一局部截面定义了 $\Sigma/\tau$ 上的连续复值
函数
\begin{equation}\label{eq:2.2}
\tilde f:\ V:=U/\tau\longrightarrow\C,\qquad [z]\longmapsto\psi\circ f(z),
\end{equation}
其中 $U$ 是 $\Sigma$ 的 $\tau$-不变开子集，$\psi$ 是折叠映射
\begin{equation}\label{eq:2.3}
\psi:\ \C\longrightarrow\C^+,\qquad x+iy\longmapsto x+i|y|
\end{equation}
（回忆我们把 $\C/\sigma$ 与闭上半平面 $\{z\in\C\mid\operatorname{Im}z\ge0\}\subset\C$
等同起来）。容易验证 (2.2) 在 $U/\tau$ 上是良定义的（因为若
$f\in(p_*\calO_\Sigma)^\tau$，则 $f(\tau(z))=f(z)$，从而
$\psi\circ f(\tau(z))=\psi\circ f(z)$）。

\begin{proposition}\label{prop:2.1.2}
给定 Real 黎曼面 $(\Sigma,\tau)$，典范投影 $p:\Sigma\to\Sigma/\tau$ 在
$\Fix(\tau)\subset\Sigma$ 与 $\Sigma/\tau$ 的边界之间诱导一个微分同胚。
$(p_*\calO_\Sigma)^\tau$ 的任一局部截面在 $\partial(\Sigma/\tau)$ 上取实值，并且
$(p_*\calO_\Sigma)^\tau$ 是 $\Sigma/\tau$ 上连续函数层的子层。
\end{proposition}

\begin{proof}
$(p_*\calO_\Sigma)^\tau$ 的局部截面是定义在 $\Sigma$ 的某个 $\tau$-不变开子集 $U$ 上
的全纯函数 $f:U\to\C$，且对一切 $x\in U$ 满足 $f\bigl(\tau(x)\bigr)=f(x)$。特别地，
若 $\tau(x)=x$，则 $f(x)\in\R$。
\end{proof}

作为 Real 黎曼面的简单例子，考虑黎曼球面 $\CP^1=\C\cup\{\infty\}$ 带实结构
$\tau:z\mapsto\bar z$，其不动点集是圆周 $\RP^1=\R\cup\{\infty\}$。$\CP^1$ 上的另一个
实结构由 $\tau_0:z\mapsto-1/\bar z$ 给出，它没有不动点。这两种实结构是 $\CP^1$ 上
仅有的可能（见习题 2.1.17）。亏格 $1$ 紧曲线上的实结构之例见习题 2.1.18。

\subsection{实曲线的拓扑型}
\label{sec:02.1.2}

本小节的目标是提供关于实代数曲线 $(\Sigma,\tau)$ 的拓扑直观，以及一个非常有用的
分类结果。这从对不动点集 $\Sigma^\tau$ 的拓扑理解开始（拓扑曲面 $\Sigma/\tau$ 在后文
也会起作用）。本节自始至终假定 $\Sigma$ 紧且连通，并记其亏格为 $g$。

\begin{proposition}\label{prop:2.1.3}
设 $(\Sigma,\tau)$ 是紧连通 Real 黎曼面。则不动点集 $\Sigma^\tau:=\Fix(\tau)$ 是
有限个圆周之并。
\end{proposition}

\begin{proof}
设 $x\in\Sigma$，$(U,\phi)$ 是 $x$ 处的 $\tau$-不变全纯坐标卡。在该卡中，映射
$z\mapsto\tau(z)$ 是双全纯映射，故局部全纯共轭于 $z\mapsto\bar z$（一般全纯映射局部
全纯共轭于 $z\mapsto z^k$，例如可参看 [18]；当且仅当 $k=1$ 时它是双射）。于是 $\tau$
局部全纯共轭于 $z\mapsto\bar z$。特别地，$\Sigma^\tau=\Fix(\tau)$ 是一维实（解析）
流形。由于 $\Sigma$ 紧，$\Sigma^\tau$ 紧且有有限个连通分支，每个分支是紧连通的一维
实流形，因此微分同胚于圆周。
\end{proof}
